\chapter{Wärme und Elektrostatik}\label{chap:scalar}
\section{Stationäre Wärme: starke und schwache Form}
Fouriers Gesetz und stationäre Energiebilanz lauten
\begin{equation}\label{eq:heat-strong}
\mathbf q=-\mathbf k\nabla T,\qquad
\nabla\cdot\mathbf q=Q,\qquad-\nabla\cdot(\mathbf k\nabla T)=Q.
\end{equation}
Auf $\Gamma_T$ ist $T=\bar T$, auf $\Gamma_q$ gilt
$\mathbf q\cdot\mathbf n=\bar q$ und auf $\Gamma_h$
$\mathbf q\cdot\mathbf n=h(T-T_\infty)$ mit $h\ge0$.
Neumann-Fluss und Konvektion können auf derselben Fläche additiv wirken.
Mit $w=0$ auf $\Gamma_T$ ergibt partielle Integration
\begin{align}\label{eq:heat-weak}
\int_\Omega\nabla w^T\mathbf k\nabla T
+\int_{\Gamma_h}whT
=\int_\Omega wQ-\int_{\Gamma_q}w\bar q
+\int_{\Gamma_h}whT_\infty.
\end{align}
Maße $\mathrm dV,\mathrm dA$ werden in dieser Gleichung zur Lesbarkeit
unterdrückt. Ein positiver äußerer Fluss \emph{entzieht} Wärme; das negative
Lastvorzeichen ist keine frei wählbare Softwareoption.
\index{Wärmeleitung}\index{Konvektion}

Aus $T_h=\mathbf N\mathbf T_e$ folgt pro Zelle
\begin{align}\label{eq:thermal-discrete}
\mathbf K_T^e&=\int_K\mathbf G^T\mathbf k\mathbf G\,\mathrm dV,&
\mathbf f_Q^e&=\int_K\mathbf N^T Q\,\mathrm dV,\\
\mathbf H^F&=\int_F h\mathbf N^T\mathbf N\,\mathrm dA,&
\mathbf f_{\infty}^F&=\int_F\mathbf N^ThT_\infty\,\mathrm dA,\nonumber\\
\mathbf f_q^F&=-\int_F\mathbf N^T\bar q\,\mathrm dA.&&\nonumber
\end{align}
Global ist $(\mathbf K_T+\mathbf H)\mathbf T=\mathbf f_Q+\mathbf f_q+
\mathbf f_\infty$. Steifigkeit und Robinmatrix tragen W/K; Lasten W.
Konstantes Tet4-Material gibt $V\mathbf G^T\mathbf k\mathbf G$.
Konstante Temperaturen liegen im Nullraum ohne Dirichlet/Robin-Bezug.
Eine positive Robinfläche kann den thermischen Nullraum entfernen.
Rein isolierte oder Neumann-Probleme benötigen kompatible Nettoleistung,
bleiben aber ohne expliziten Bezug uneindeutig; die API wählt keinen Mittelwert.

\subsection{Leistungsbilanz und Ergebnisfelder}
Für die diskrete Einsfunktion fallen innere Leitbeiträge weg. Mit der
Dirichlet-Reaktion als Zuflussleistung folgt
\begin{equation}\label{eq:thermal-balance}
P_Q+P_{D,\rm in}-P_{q,\rm out}-P_{h,\rm out}=0.
\end{equation}
$P_{D,\rm in}=\sum r_i$ über feste Temperatur-DOFs, während
$P_{q,\rm out}=-\sum f_{q,i}$. Robinabfluss wird aus
$\mathbf H\mathbf T-\mathbf f_\infty$ integriert.
Die Bilanz verwendet Lasten/Reaktionen, nicht eine unbegründete Integration
von unstetigen P1-Flussproben. Tet4-Flüsse sind elementkonstant; Tet10-Flüsse
werden als Schwerpunktproben mit Auswerteort berichtet.

\subsection{Handrechnung: Leitung und Konvektion}
Ein 1-m-Quader mit $k=2$ W/(m K), $T(0)=400$ K, $Q=0$ und
$h=4$ W/(m$^2$ K), $T_\infty=300$ K bei $x=1$ besitzt konstanten Fluss.
Aus $T(x)=400-qx/k$ und $q=h(T(1)-300)$ folgt
\begin{equation}\label{eq:thermal-example}
q=\frac{100}{1/2+1/4}=\frac{400}{3}\ {\rm W/m^2},
\qquad T(1)=\frac{1000}{3}\ {\rm K}.
\end{equation}
Bei Fläche 1 m$^2$ sind Zufluss und Konvektionsabfluss je $400/3$ W.
Mit vorgeschriebenem Fluss 20 W/m$^2$ statt Robin gilt $T=400-10x$ K.
Beide affinen Felder sind auf konformen Tet4/Tet10-Netzen exakt darstellbar.

\section{Elektrostatik: Gauß-Gesetz und Vorzeichen}
Für ein lineares Dielektrikum gilt
\begin{equation}\label{eq:electro-strong}
\mathbf E=-\nabla\phi,\qquad\mathbf D=\boldsymbol\kappa\mathbf E,
\qquad\nabla\cdot\mathbf D=\rho_f.
\end{equation}
Potential ist auf $\Gamma_\phi$ vorgegeben, auf $\Gamma_\omega$
$\mathbf D\cdot\mathbf n=\bar\omega$. Multiplizieren mit $\psi$, integrieren
und Umstellen ergibt die Projektform
\begin{equation}\label{eq:electro-weak}
\int_\Omega\nabla\psi^T\mathbf D\,\mathrm dV
=\int_{\Gamma_\omega}\psi\bar\omega\,\mathrm dA-
\int_\Omega\psi\rho_f\,\mathrm dV.
\end{equation}
Daher ist $-\mathbf K_\phi\boldsymbol\phi=\mathbf f_\phi$ mit
\begin{equation}\label{eq:dielectric}
\mathbf K_\phi^e=\int_K\mathbf G^T\boldsymbol\kappa\mathbf G\,\mathrm dV,
\quad\mathbf f_\phi=\mathbf f_\omega-\mathbf f_\rho,
\quad\mathbf f_\rho^e=\int_K\mathbf N^T\rho_f\,\mathrm dV.
\end{equation}
$\mathbf K_\phi$ ist positiv semidefinit und hat Einheit F; das Minus steht
im assemblierten Operator. Der Einzelfeldsolver multipliziert beide freien
Gleichungsseiten gemeinsam mit $-1$ für ein positiv definites Solve-System.
Reaktionen stammen weiterhin aus der ursprünglichen negativen Gleichung.
Dieses Vorgehen erhält die spätere Piezo-Symmetrie.
\index{Elektrostatik}\index{Gauge}

Jede zusammenhängende Komponente braucht einen Potentialbezug. Für reine
Ladungsdaten ist $\int\rho_f=\int\bar\omega$ nötig, aber die additive Konstante
bleibt frei. Der Solver diagnostiziert kompatible bzw. inkompatible Daten
und erdet nicht automatisch. Ein Potentialoffset verändert $\mathbf E$ nicht.

\section{Ladung, Leiter und positive Feldenergie}
Die Ladung des angrenzenden feldfreien Leiters ist das Negative des
D-Flusses mit der \emph{Dielektrikumsnormalen}. Auf einem Potentialrand ist
$Q_L=-\sum_i r_i$. Die diskrete Bilanz lautet
\begin{equation}\label{eq:charge-balance}
Q_{\omega,\rm prescribed}+Q_{D,\rm reaction}-Q_{\rho,\rm volume}=0.
\end{equation}
Die positive gespeicherte Dielektrikumsenergie ist
\begin{equation}\label{eq:electric-energy}
U_E=\tfrac12\int_\Omega\mathbf E^T\boldsymbol\kappa\mathbf E\,\mathrm dV
=\tfrac12\boldsymbol\phi^T\mathbf K_\phi\boldsymbol\phi\ge0.
\end{equation}
Ohne Volumen-/vorgegebene Randladung ist $2U_E=\sum\phi_i Q_{L,i}$.
Der negative Enthalpiebeitrag in der Piezoform darf damit nicht verwechselt
werden (Kapitel~\ref{chap:coupling}).

\subsection{Handrechnung: idealer Plattenkondensator}
$L=0{,}002$ m, $A=0{,}03\cdot0{,}04=0{,}0012$ m$^2$,
$\kappa=2\cdot10^{-11}$ F/m, $V=100$ V. Seiten sind D-flussfrei;
Außenraum/Randfelder werden nicht modelliert. Dann
\begin{equation}\label{eq:capacitor-example}
C_{\rm cap}=\kappa A/L=12\ {\rm pF},\quad
Q_{L,\rm high}=1{,}2\ {\rm nC},\quad U_E=60\ {\rm nJ}.
\end{equation}
$E_x=-50000$ V/m. Am linken Dielektrikumsrand ist $D\cdot n>0$,
seine Leiterladung negativ; an der Hochpotentialplatte ist sie positiv.
Zwei gleich dicke Schichten in Serie erfüllen
$C_{\rm cap}^{-1}=L_1/(\kappa_1A)+L_2/(\kappa_2A)$.
Das kontinuierliche D-Feld darf deshalb nicht durch Gleichheit der
Gradienten an einer Materialgrenze ersetzt werden.

\subsection*{Implementierung und unabhängige Nachweise.}
{\small\raggedright
\path{src/fem_toolbox/physics/heat.py},
\path{src/fem_toolbox/physics/electrostatics.py},
\path{src/fem_toolbox/solvers/heat.py},
\path{src/fem_toolbox/solvers/electrostatics.py};
\path{tests/patch/test_heat_patch.py},
\path{tests/patch/test_electrostatics_patch.py},
\path{tests/patch/test_material_interfaces.py}.
Vier-Netze-Nachweise: \path{tests/convergence/test_heat_convergence.py},
\path{tests/convergence/test_electrostatics_convergence.py} und die
entsprechenden Tet10-Tests. Doc/11, Doc/12 und Doc/45--46 legen API und
maßstabsabhängige absolute Toleranzen fest.

\par}
